\chapter{Problem Setting}
\label{ch:problem}

This chapter describes the world the robot works in, the choices it has, and the yardstick by
which its behaviour is judged. Everything stated here is a property of the problem, fixed before
any solution is designed; the decision models that answer it follow in \cref{ch:models}. The
chapter closes by showing, with the numbers of the setting itself, why the question
\emph{what should the robot do next?} has no obvious answer.

\section{The factory}
\label{sec:factory}

The setting is a single production hall of \SI{14.5}{\metre} by \SI{10}{\metre}, laid out as a
maze of walls, machines and corridors (\cref{fig:maze}). Three production sections, named A, B and
C, stand in different regions of the hall. Each of them manufactures finished units and places
them in a local output buffer. A single delivery point in the south of the hall accepts finished
units, and a single charging station in the south-west corner is where the robot starts its shift
and the only place it can recharge.

\begin{figure}[H]
    \centering
    \includegraphics[width=0.92\textwidth]{maze_layout.png}
    \caption{The production hall: three sections (A north-west, B north-east, C south-east), one
             delivery point and one charging station. Every corridor is at least
             \SI{1.5}{\metre} wide, and the layout is deliberately asymmetric so that different
             places produce distinguishable sensor readings.}
    \label{fig:maze}
\end{figure}

The geometry is not decorative. Two properties of it shape the decision problem. First, the
distances between the points of interest differ by almost a factor of three
(\cref{tab:distances}), so the order in which sections are served changes how far the robot
drives and therefore how much energy it spends. Second, several pairs of points are connected by
more than one route, and the shortest route is not always the one a planner finds --- a fact that
becomes important in \cref{sec:planner-study}.

\begin{table}[H]
    \centering
    \caption{Shortest path through the maze between the points of interest, in metres. These are
             the true walkable distances, computed on the occupancy grid, not straight lines.}
    \label{tab:distances}
    \begin{tabular}{lrrrrr}
        \toprule
                 & A & B & C & delivery & charger\\
        \midrule
        A        & --- & 18.28 & 15.56 & 10.01 & 9.04\\
        B        & 18.28 & --- & 9.24 & 10.26 & 15.18\\
        C        & 15.56 & 9.24 & --- & 6.30 & 11.27\\
        delivery & 10.01 & 10.26 & 6.30 & --- & 5.10\\
        charger  & 9.04 & 15.18 & 11.27 & 5.10 & ---\\
        \bottomrule
    \end{tabular}
\end{table}

The three sections are deliberately unequal (\cref{tab:sections}). Section~B is the fast line: it
produces one unit every two minutes, and its buffer of ten units is full after twenty minutes of
neglect. Section~A is half as fast with a slightly larger reserve of time. Section~C produces only
one unit every twelve minutes, but its units are almost four times as heavy as B's and its buffer holds
just three of them, so it can be ignored for a while --- but not forgotten. Production is not
perfectly regular: cycle times are drawn from a Gamma distribution with a coefficient of variation
of \num{0.15}, so buffers do not fill on a predictable schedule.

\begin{table}[H]
    \centering
    \caption{The three production sections at nominal rates. ``Time to full'' assumes an empty
             buffer and no interruption.}
    \label{tab:sections}
    \begin{tabular}{lrrrrr}
        \toprule
        Section & Rate (units/h) & Buffer (units) & Unit mass (kg) & Full buffer (kg) & Time to full\\
        \midrule
        A & 20 & 8  & 0.5 & 4.0 & \SI{24}{\minute}\\
        B & 30 & 10 & 0.4 & 4.0 & \SI{20}{\minute}\\
        C & 5  & 3  & 1.5 & 4.5 & \SI{36}{\minute}\\
        \bottomrule
    \end{tabular}
\end{table}

What turns this into a problem rather than a chore is the consequence of a full buffer. A section
whose buffer is full cannot put the next finished unit anywhere, so the line stops: it enters the
state \textsc{blocked}, and everything it would have produced while blocked is lost for good. Lost
production cannot be recovered later by working harder, which is precisely why the robot's
decisions matter. Two further states model the wear of the machines themselves. A machine's health
decreases slowly with every unit it makes; below \SI{60}{\percent} health it runs at a reduced
rate (\textsc{degraded}), and it can break down entirely (\textsc{fault}), which stops production
for a repair of about twenty minutes. Faults are more likely the less healthy a machine is.

\section{The robot}
\label{sec:robot}

The transport robot is a small differential-drive platform of the TurtleBot class: \SI{4.3}{\kilo\gram},
a two-dimensional lidar for localisation, and a single-board computer. Its measured average speed
through this maze, turns included, is \SI{0.39}{\metre\per\second}. It carries at most
\SI{5}{\kilo\gram} at a time, and loading or unloading takes \SI{30}{\second} of standing still.

Four physical properties constrain what it can do, and all four are modelled explicitly rather
than assumed away (the full parameter set is in \cref{app:configuration}):

\begin{description}[leftmargin=1.2em, style=nextline]
  \item[Energy.] The robot draws \SI{6}{\watt} whenever it is switched on, plus
        \SI{0.005}{\watt\hour\per\metre} while driving empty and a further
        \SI{0.00116}{\watt\hour\per\metre} for every kilogram it carries. Its battery holds
        \SI{22}{\watt\hour}, which is roughly \SI{1.7}{\hour} of continuous driving --- enough for
        a shift only if the robot does not drive needlessly.
  \item[Charging.] The charging station delivers \SI{25}{\watt}, so a full recharge takes about an
        hour. Charging is therefore not a pause but a decision with a real cost in lost working
        time, which is why the robot can choose \emph{how far} to charge rather than only whether
        to.
  \item[Heat.] The motors warm up under load and cool down when idle, reaching a steady state of
        about \SI{55}{\degreeCelsius} at continuous full load in a \SI{22}{\degreeCelsius} hall.
        Above \SI{55}{\degreeCelsius} wear accelerates; at \SI{70}{\degreeCelsius} the robot must
        stop and cool down. A heavy payload heats it faster.
  \item[Wear.] The drivetrain's condition degrades with energy drawn and with time spent hot. A
        worn drive is not merely old: at \SI{0}{\percent} condition it draws \SI{30}{\percent} more
        energy for the same journey, and below \SI{30}{\percent} the robot is due for maintenance
        and may no longer be given work.
\end{description}

These properties interact. Carrying more units per trip saves journeys but costs energy per metre
and heats the motors faster; driving with a hot motor costs extra energy through reduced
efficiency; and every watt-hour drawn wears the drive a little further, which in turn makes future
journeys more expensive.

\section{What the robot can do}
\label{sec:actions}

The robot's behaviour is expressed as a sequence of complete actions rather than as continuous
motion. One action is one decision, and the next decision is only taken when the current action
has finished. Four kinds of action exist:

\begin{description}[leftmargin=1.2em, style=nextline]
  \item[\texttt{PICKUP:<A|B|C>}] Drive to the named section and load as many waiting units as
        still fit within the payload limit, emptying that much of its buffer.
  \item[\texttt{DELIVER}] Drive to the delivery point and unload everything carried. Only delivered
        units count as produced value.
  \item[\texttt{CHARGE:<40|70|100>}] Drive to the charging station and charge up to the chosen
        level. Offering three levels makes the trade-off explicit: a quick top-up returns the robot
        to work sooner, a full charge buys independence for longer.
  \item[\texttt{WAIT:<s>}] Stand still for sixty seconds. At the charging station this is free and
        tops up the battery; anywhere else it merely spends idle power.
\end{description}

Driving itself is not part of the decision problem. A navigation stack plans and follows the route,
recovers from obstacles, and reports success or failure (\cref{sec:executor}); the decision layer
chooses \emph{where} to go, never \emph{how} to get there.

\section{What counts as a good shift}
\label{sec:objective}

A shift lasts one hour of factory time. Everything the robot does during it is judged by a single
number, so that different behaviours can be compared without argument:

\begin{equation}
    \label{eq:objective}
    J \;=\; n_{\text{delivered}}
        \;-\; n_{\text{lost}}
        \;-\; 0.5\,E_{\text{battery}}
        \;-\; 20\,n_{\text{violations}}
\end{equation}

where $n_{\text{delivered}}$ is the number of units brought to the delivery point,
$n_{\text{lost}}$ the production lost while lines stood blocked, $E_{\text{battery}}$ the energy
drawn from the battery in watt-hours, and $n_{\text{violations}}$ the number of safety violations
(defined in \cref{sec:constraints}).

The weights encode a position that is worth stating plainly. A delivered unit and a lost unit
weigh the same, because a unit that never got made is as costly as one that never arrived. Energy
is deliberately cheap: at \num{0.5} per watt-hour, and with a measured journey between two points of the hall costing
roughly \SIrange{0.18}{0.25}{\watt\hour}, no reasonable policy can improve its score by refusing to work --- energy is a
tie-breaker between journeys, not a reason to stay at home. A safety violation, in contrast, costs
twenty units, far more than any shift's worth of careful energy saving, because a stranded or
overheated robot is not a slightly worse outcome but a different kind of outcome: it requires a
human being to walk into the hall.

Two things this objective deliberately does not charge for. Time spent idle is free, since the
shift has a fixed length and an idle robot in a calm factory harms nothing. And only energy drawn
\emph{from the battery} is counted, not the energy the charging station supplies while the robot
is docked. The second choice is a simplification whose consequences are revisited honestly in
\cref{sec:limitations}.

\section{Constraints}
\label{sec:constraints}

Four conditions must hold at all times, and breaking any of them counts as a safety violation in
\cref{eq:objective}:

\begin{enumerate}[label=\textbf{C\arabic*.}, leftmargin=*]
  \item \textbf{Battery reserve.} At least \SI{15}{\percent} of the pack must remain after the robot
        has returned to the charging station. The reserve is not a threshold on the current battery
        level but on the level \emph{after the journey home}, which means the constraint depends on
        where the robot is and what it is carrying.
  \item \textbf{Temperature.} Motor temperature must stay below \SI{70}{\degreeCelsius}.
  \item \textbf{Payload.} The robot may never plan to carry more than \SI{5}{\kilo\gram}.
  \item \textbf{Maintenance.} Below \SI{30}{\percent} drive condition the robot may only charge or
        wait; it is no longer fit for transport work.
\end{enumerate}

Constraint~C1 is the interesting one, because it cannot be checked by looking at the battery
gauge. Answering \emph{may I fetch from section~B?} requires simulating the whole journey ---
there, loaded, and home again --- through the robot's own energy, thermal and wear models. The
same question asked with a hot motor, a heavy load or a worn drive has a different answer.

\section{Why the problem is not trivial}
\label{sec:why-hard}

The setting was designed so that no simple rule is obviously right. Three pieces of arithmetic
show why.

\paragraph{The robot is not comfortably faster than the factory.}
At nominal rates the three lines together produce \num{55} units per hour with a combined mass of
\SI{29.5}{\kilo\gram}. With a payload limit of \SI{5}{\kilo\gram}, at least six loaded deliveries
per hour are needed just to keep pace, and each of them costs travel time to the section, loading,
travel to the delivery point and unloading. The robot has enough capacity for the nominal factory,
but not much to spare --- and the scenarios in \cref{sec:scenarios} include days where the demand
exceeds what one robot can carry, so that some loss is unavoidable and the only question is how
little.

\paragraph{Urgency and cost point in different directions.}
Section~B fills fastest and is therefore usually the most urgent, but it is also the section
farthest from the charging station (\SI{15.18}{\metre}) and from section~A
(\SI{18.28}{\metre}). Section~C fills slowest, but sits closest to the delivery point
(\SI{6.30}{\metre}) and carries the heaviest units, so serving it is cheap per journey and
expensive per kilogram. A policy that always serves the most urgent line drives further than one
that batches nearby work; a policy that always minimises distance eventually lets a line block.

\paragraph{Charging competes with everything else.}
Recharging from empty takes about an hour at \SI{25}{\watt} --- a whole shift. A robot that charges
too early gives up working time it cannot get back; one that charges too late violates C1, and the
penalty for that is worth twenty delivered units. Because the reserve must hold \emph{after} the
journey home, the moment at which charging becomes unavoidable depends on where the robot happens
to be standing.

None of these tensions can be resolved once and written down as a fixed order of visits, because
the state that decides them --- buffer fills, machine health, battery, temperature, condition ---
changes continuously and partly at random.

\section{Scenarios}
\label{sec:scenarios}

A single factory configuration would measure very little, so the problem is instantiated in twelve
scenarios (\cref{tab:scenarios}). Each one changes the parameters of \cref{sec:factory} or
\cref{sec:robot} --- production rates, buffer sizes, unit masses, machine health, ambient
temperature, battery capacity or the robot's initial state --- while the objective and the
constraints stay the same.

Six of the scenarios were written before the decision models were trained, and both learned models
in \cref{ch:models} use exactly those six. The other six were written afterwards, when the models
already existed, and are never used for training or for selecting a model. They therefore measure
something the first six cannot: whether a policy still behaves sensibly in a factory it has never
seen. This separation is what makes \cref{sec:generalisation} a measurement rather than a claim,
and it is maintained down to the random seeds (\cref{sec:splits}).

\begin{table}[H]
    \centering
    \caption{The twelve scenarios. ``Seen'' scenarios are available for training; ``unseen''
             scenarios were written after training and are used only for evaluation.}
    \label{tab:scenarios}
    \begin{tabularx}{\textwidth}{llX}
        \toprule
        Scenario & Split & What it stresses\\
        \midrule
        \texttt{balanced}          & seen   & Nominal rates, rare faults, healthy fully charged robot.\\
        \texttt{high\_demand}      & seen   & All lines at twice the nominal rate; robot utilisation about \num{0.85}.\\
        \texttt{one\_hot\_section} & seen   & Section~B at double speed, blocking within ten minutes if neglected.\\
        \texttt{fault\_burst}      & seen   & Worn machines failing about once per hour with long repairs.\\
        \texttt{low\_battery\_start} & seen & Robot starts at \SI{30}{\percent} battery with buffers already half full.\\
        \texttt{worn\_robot}       & seen   & Drive condition \SI{45}{\percent}: more energy per metre, close to the maintenance limit.\\
        \midrule
        \texttt{heavy\_load}       & unseen & Double production, heavier parts, half-charged robot: slightly beyond one robot's capacity.\\
        \texttt{heavy\_parts}      & unseen & Units about twice as heavy, so fewer fit per trip.\\
        \texttt{small\_buffers}    & unseen & Lean production: buffers hold only a few units and block within minutes.\\
        \texttt{section\_breakdown}& unseen & Section~B's machine failing (health \SI{35}{\percent}) while A and C run normally.\\
        \texttt{hot\_factory}      & unseen & \SI{38}{\degreeCelsius} in the hall and poor cooling: long heavy trips risk the temperature limit.\\
        \texttt{aged\_battery}     & unseen & Pack down to \SI{12}{\watt\hour} (\SI{55}{\percent} of new) and a slower \SI{15}{\watt} charger.\\
        \bottomrule
    \end{tabularx}
\end{table}

With the world, the actions, the objective and the constraints fixed, the question of this thesis
can be stated exactly: \emph{which policy, choosing one action at a time from the state of the
factory and the robot, maximises \cref{eq:objective} without ever breaking C1--C4 --- including in
factories it was never designed or trained for?}
